\section{Problem 3: Sensor Calibration}

\subsection{Problem description}
The objective is to determine the Steinhart-Hart coefficients (A, B, C) that model the relationship between Resistance (R) and Temperature (T) for a 10K thermistor.
Model Equation: 
\begin{equation}
T_{model} = \frac{1}{A + B \ln(R) + C (\ln(R))^3}
\end{equation}
Note: $T$ is in Kelvin, $R$ in Ohms.

\subsection{Cost function}
The optimisation goal is to minimise the Sum of Squared Errors (SSE) between the measured temperature and the model-predicted temperature:
\begin{equation}
SSE = \sum (T_{measured,i} - T_{model,i})^2
\end{equation}

\subsection{Sources}
\begin{itemize}
    \item \url{https://www.bapihvac.com/wp-content/uploads/2010/11/Thermistor_10K-2.pdf}
    \item \url{https://www.sciencedirect.com/science/article/abs/pii/0011747168900570}
\end{itemize}

\subsection{Methodology}
\textbf{Simulated Annealing:} A probabilistic technique for approximating the global optimum of a given function.

\textbf{Particle Swarm Optimisation (PSO):} A computational method that optimises a problem by iteratively trying to improve a candidate solution based on a given measure of quality.

\subsection{Settings}

\textbf{Settings SA:}
\begin{itemize}
    \item Initial Temperature: 50.0
    \item Cooling Rate: 0.995
    \item Minimum Temperature: 1e-6
\end{itemize}

\textbf{Setting PSO:} Implementation: Utilises the pymoo library (PSO algorithm).
\begin{itemize}
    \item pop\_size: Population size (controls diversity).
    \item w: Inertia weight (balance between exploration and exploitation).
    \item c1: Cognitive parameter (individual experience).
    \item c2: Social parameter (swarm experience).
    \item max\_velocity\_rate: Limits the maximum velocity of particles.
\end{itemize}

\subsection{Result}

The optimal Steinhart-Hart coefficients found using Simulated Annealing for the provided dataset are:

\begin{itemize}
    \item A = $1.114203 \times 10^{-3}$
    \item B = $2.366744 \times 10^{-4}$
    \item C = $7.761151 \times 10^{-8}$
\end{itemize}

Final SSE: 0.4323

\begin{figure}[h!]
    \centering
    \includegraphics[width=0.8\textwidth]{best_solution_model.png}
    \caption{Steinhart-Hart Model Fit}
\end{figure}

\subsection{Work discussion}
The coefficients obtained result in a very low Sum of Squared Errors (0.4323), indicating that the Steinhart-Hart model fits the thermistor data accurately. The use of Simulated Annealing allowed for robust convergence even with random initialization.
